% ============================================================
%  MANUAL DE USUARIO — AI Presentation Analyzer
%  Trabajo Terminal II · 2026-B127 · ESCOM — IPN
%
%  Cómo compilar:  tectonic "Manual de Usuario.tex"
%  (también funciona con pdflatex, dos pasadas)
% ============================================================
\documentclass[11pt]{article}

% ---------- Página ----------
\usepackage[letterpaper,margin=2.2cm]{geometry}
\usepackage{iftex}
\ifxetex
  \usepackage{fontspec}
  \setmainfont{lmroman10-regular.otf}[BoldFont=lmroman10-bold.otf,ItalicFont=lmroman10-italic.otf,BoldItalicFont=lmroman10-bolditalic.otf]
  \setsansfont{lmsans10-regular.otf}[BoldFont=lmsans10-bold.otf]
  \setmonofont{lmmono10-regular.otf}
\else
  \usepackage[utf8]{inputenc}
  \usepackage[T1]{fontenc}
  \usepackage{lmodern}
\fi
\usepackage[spanish,es-noshorthands]{babel}

% ---------- Gráficos y color ----------
\usepackage[table]{xcolor}
\usepackage{graphicx}
\usepackage{tikz}
\usetikzlibrary{arrows.meta}
\usepackage{float}

% ---------- Tablas y listas ----------
\usepackage{array}
\usepackage{booktabs}
\usepackage{tabularx}
\usepackage{enumitem}

% ---------- Cajas y estilos ----------
\usepackage{mdframed}
\usepackage{titlesec}
\usepackage{fancyhdr}

% ---------- Enlaces ----------
\PassOptionsToPackage{hyphens}{url}
\usepackage{url}
\usepackage[colorlinks=true]{hyperref}

% ============================================================
%  PALETA (guinda institucional + grises neutros)
% ============================================================
\definecolor{guinda}{HTML}{6A1B31}
\definecolor{guindaoscuro}{HTML}{4D1324}
\definecolor{guindasuave}{HTML}{FDF2F4}
\definecolor{oro}{HTML}{BC955C}
\definecolor{gristexto}{HTML}{475569}
\definecolor{grismedio}{HTML}{94A3B8}
\definecolor{grisclaro}{HTML}{E2E8F0}
\definecolor{grisfondo}{HTML}{F8FAFC}
\definecolor{pizarron}{HTML}{334155}

\hypersetup{
  linkcolor=gristexto,
  urlcolor=gristexto,
  citecolor=gristexto,
  pdftitle={Manual de Usuario — AI Presentation Analyzer},
  pdfauthor={Escobar Segura Brandon Aldair, Gonzalez Ruiz Karol},
  pdfsubject={Manual de usuario del prototipo de análisis de presentaciones impulsado por IA},
}

% ============================================================
%  FORMATO GENERAL
% ============================================================
\setlength{\parindent}{0pt}
\setlength{\parskip}{5pt}
\setlength{\emergencystretch}{3em}
\renewcommand{\arraystretch}{1.4}
\arrayrulecolor{grisclaro}

\titleformat{\section}
  {\sffamily\LARGE\bfseries\color{pizarron}}{\textcolor{guinda}{\thesection}}{0.55em}{}
  [{\color{grisclaro}\titlerule[1pt]\vspace{1pt}}]
\titlespacing*{\section}{0pt}{26pt}{14pt}
\titleformat{\subsection}
  {\sffamily\large\bfseries\color{pizarron}}{\textcolor{guinda}{\thesubsection}}{0.55em}{}
\titlespacing*{\subsection}{0pt}{16pt}{8pt}

\pagestyle{fancy}
\fancyhf{}
\setlength{\headheight}{14.5pt}
\fancyhead[L]{\sffamily\scriptsize\textcolor{gristexto}{Manual de Usuario — AI Presentation Analyzer}}
\fancyhead[R]{\sffamily\scriptsize\textcolor{guinda}{\textbf{ESCOM · IPN}}}
\fancyfoot[L]{\sffamily\scriptsize\textcolor{gristexto}{Trabajo Terminal 2026-B127}}
\fancyfoot[R]{\sffamily\scriptsize\textcolor{gristexto}{\textbf{Página \thepage}}}
\renewcommand{\headrulewidth}{0pt}
\renewcommand{\headrule}{{\color{grisclaro}\hrule width\headwidth height 0.6pt}}

\addto\captionsspanish{%
  \renewcommand{\tablename}{Tabla}%
  \renewcommand{\contentsname}{\sffamily\color{guinda}Contenido}%
}

% ============================================================
%  ELEMENTOS DE DISEÑO
% ============================================================

% --- Botón de la interfaz (relleno guinda, texto blanco) ---
\newcommand{\btn}[1]{{\setlength{\fboxsep}{3.5pt}\setlength{\fboxrule}{0pt}%
\fcolorbox{guinda}{guinda}{\textcolor{white}{\sffamily\bfseries\footnotesize #1}}}}

% --- Botón suave (fondo guinda pálido, texto guinda) ---
\newcommand{\btnsuave}[1]{{\setlength{\fboxsep}{3.5pt}\setlength{\fboxrule}{0pt}%
\fcolorbox{guindasuave}{guindasuave}{\textcolor{guinda}{\sffamily\bfseries\footnotesize #1}}}}

% --- Elemento de pantalla (fondo gris claro) ---
\newcommand{\ui}[1]{{\setlength{\fboxsep}{3pt}\setlength{\fboxrule}{0.6pt}%
\fcolorbox{grisclaro}{grisfondo}{\textcolor{gristexto}{\sffamily\footnotesize #1}}}}

% --- Tecla / atajo ---
\newcommand{\tecla}[1]{{\setlength{\fboxsep}{2.5pt}\setlength{\fboxrule}{0.6pt}%
\fcolorbox{grisclaro}{white}{\textcolor{pizarron}{\ttfamily\footnotesize\bfseries #1}}}}

% --- Número de paso (círculo gris oscuro) ---
\newcommand{\numeropaso}[1]{\raisebox{-2pt}{\tikz[baseline=(n.base)]%
\node[circle,fill=pizarron,text=white,font=\sffamily\bfseries\scriptsize,%
inner sep=1pt,minimum size=15.5pt](n){#1};}}

\newlist{pasos}{enumerate}{1}
\setlist[pasos]{label=\protect\numeropaso{\arabic*},leftmargin=27pt,labelsep=9pt,%
itemsep=8pt,topsep=8pt}

\setlist[itemize,1]{label=\textcolor{grismedio}{\textbullet},leftmargin=18pt,itemsep=3pt,topsep=4pt}

% --- Caja neutral (notas) ---
\newmdenv[backgroundcolor=grisfondo,linecolor=grisclaro,linewidth=1pt,%
innerleftmargin=10pt,innerrightmargin=10pt,innertopmargin=9pt,innerbottommargin=9pt,%
skipabove=12pt,skipbelow=12pt]{cajanotaenv}
\newenvironment{nota}[1][Nota]
  {\begin{cajanotaenv}{\sffamily\bfseries\color{pizarron}#1}\par\vspace{2pt}}
  {\end{cajanotaenv}}

% --- Caja de advertencia (barra guinda a la izquierda) ---
\newmdenv[backgroundcolor=grisfondo,linecolor=guinda,linewidth=3.5pt,%
innerleftmargin=10pt,innerrightmargin=10pt,innertopmargin=9pt,innerbottommargin=9pt,%
skipabove=12pt,skipbelow=12pt]{cajaimportanteenv}
\newenvironment{importante}[1][Importante]
  {\begin{cajaimportanteenv}{\sffamily\bfseries\color{guinda}#1}\par\vspace{2pt}}
  {\end{cajaimportanteenv}}

% --- Columnas de tabla ---
\newcolumntype{R}[1]{>{\raggedright\arraybackslash}p{#1}}
\newcolumntype{C}[1]{>{\centering\arraybackslash}p{#1}}
\newcommand{\encabezado}[1]{\textcolor{white}{\sffamily\bfseries #1}}

% ============================================================
%  CAPTURAS DE PANTALLA
%  \imagen[opciones]{archivo}  — incluye img/<archivo>.png
%  \pendiente[altura]{archivo}{descripción} — marco para captura
%  faltante (páginas web que requieren la cuenta del usuario).
% ============================================================
\newcommand{\imagen}[2][width=0.9\linewidth]{%
\begin{center}
\IfFileExists{img/#2.png}{\fcolorbox{grisclaro}{white}{\includegraphics[#1]{img/#2.png}}}%
{\fcolorbox{oro}{grisfondo}{\parbox{0.7\linewidth}{\centering
\sffamily\footnotesize\textcolor{guinda}{CAPTURA PENDIENTE}\\[2pt]
\ttfamily\scriptsize\textcolor{gristexto}{img/#2.png}}}}
\end{center}}

\newcommand{\pendiente}[3][6cm]{%
\begin{center}
\begin{tikzpicture}
\draw[color=oro,line width=1pt,dashed,rounded corners=7pt,fill=grisfondo]
  (0,0) rectangle (\linewidth,#1);
\node[align=center,text width=0.84\linewidth] at (current bounding box.center) {%
  {\sffamily\bfseries\footnotesize\textcolor{guinda}{CAPTURA DE PANTALLA PENDIENTE}}\\[4pt]
  {\sffamily\small\textcolor{gristexto}{#3}}\\[5pt]
  {\ttfamily\scriptsize\textcolor{oro}{img/#2.png}}
};
\end{tikzpicture}
\end{center}}

% ============================================================
\begin{document}

% ============================================================
%  PORTADA
% ============================================================
\pagenumbering{gobble}
\begin{titlepage}
  \thispagestyle{empty}

  % ===== ENCABEZADO INSTITUCIONAL =====
  \noindent
  \begin{minipage}[c]{0.18\textwidth}
      \centering
      \includegraphics[width=2.0cm,keepaspectratio]{IPN.png}
  \end{minipage}\hfill
  \begin{minipage}[c]{0.60\textwidth}
      \centering
      {\large\textbf{INSTITUTO POLITÉCNICO NACIONAL}}\\[0.15cm]
      {\large\textbf{ESCUELA SUPERIOR DE CÓMPUTO}}
  \end{minipage}\hfill
  \begin{minipage}[c]{0.18\textwidth}
      \centering
      \includegraphics[width=2.8cm,keepaspectratio]{ESCOM.png}
  \end{minipage}

  \vspace{1.6cm}
  \begin{center}
      {\color{guinda}\rule{\linewidth}{2.5pt}}
  \end{center}

  \vspace{1.0cm}

  % ===== TÍTULO DEL PROYECTO =====
  \begin{center}
      {\sffamily\small\bfseries\textcolor{gristexto}{Trabajo Terminal II \quad\textbullet\quad Número: 2026-B127}}\\[0.55cm]
      {\color{guinda}\LARGE\bfseries\sffamily Manual de Usuario}\\[0.5cm]
      {\sffamily\large\textcolor{gristexto}{AI Presentation Analyzer}}\\[0.35cm]
      {\sffamily\normalsize\textcolor{gristexto}{Prototipo de análisis de presentaciones\\impulsado por inteligencia artificial}}
  \end{center}

  \vspace{1.4cm}
  \begin{center}
      {\color{oro}\rule{0.35\linewidth}{1.5pt}}
  \end{center}
  \vspace{1.0cm}

  % ===== AUTORES Y DIRECTORES =====
  \noindent
  \begin{minipage}[t]{0.48\textwidth}
      \centering
      {\large\textbf{Presentan:}}\\[0.35cm]
      {\large Escobar Segura Brandon Aldair}\\[0.15cm]
      {\large Gonzalez Ruiz Karol}
  \end{minipage}\hfill
  \begin{minipage}[t]{0.48\textwidth}
      \centering
      {\large\textbf{Directores de Trabajo Terminal:}}\\[0.35cm]
      {\large Camacho Vázquez Vanessa Alejandra}\\[0.15cm]
      {\large Rosas Carrillo Ary Shared}
  \end{minipage}

  \vfill

  % ===== PIE =====
  \begin{center}
      {\large Versión 1.0 \quad\textcolor{oro}{\textbullet}\quad
       Ciudad de México \quad\textcolor{oro}{\textbullet}\quad Diciembre 2026}
  \end{center}
\end{titlepage}

% ============================================================
%  CONTENIDO
% ============================================================
\pagenumbering{roman}
\pdfbookmark[1]{Contenido}{contenido}
\tableofcontents
\newpage

\pagenumbering{arabic}

% ============================================================
%  1. LA APLICACIÓN
% ============================================================
\section{La aplicación}
\label{sec:overview}

\subsection{¿Qué es y para qué sirve?}
\label{subsec:que-es}

\textbf{AI Presentation Analyzer} es un prototipo de escritorio desarrollado como
Trabajo Terminal en la Escuela Superior de Cómputo (ESCOM) del Instituto Politécnico
Nacional. Está pensado para profesores que exponen con diapositivas en formato
\texttt{.pptx}: evalúa la calidad de sus presentaciones mediante métricas objetivas y
la inteligencia artificial de Google Gemini, entrega recomendaciones concretas de
mejora, genera una versión optimizada de cada presentación y acompaña la exposición
frente al grupo con pantallas de apoyo y un informe de rendimiento.

\subsection{¿Cómo funciona?}
\label{subsec:como-funciona}

El prototipo organiza el trabajo del profesor en cuatro momentos:

\begin{enumerate}
  \item \textbf{Preparar.} Registra las materias que imparte y carga en cada una
        sus presentaciones \texttt{.pptx}.
  \item \textbf{Analizar y mejorar.} La inteligencia artificial revisa cada
        diapositiva, muestra resultados con etiquetas de calidad, sugiere cambios de
        contenido y genera una versión optimizada que se guarda en la carpeta
        \ui{Documentos} del equipo.
  \item \textbf{Exponer.} Durante la clase, proyecta la presentación al grupo y
        conserva en su pantalla un panel de apoyo con cronómetros, preguntas y
        datos complementarios de cada diapositiva.
  \item \textbf{Dar seguimiento.} Consulta el informe de rendimiento de cada sesión
        y el historial de versiones de cada presentación.
\end{enumerate}

\subsection{Funciones que ofrece}
\label{subsec:funciones}

\begin{center}
\small
\begin{tabularx}{\linewidth}{>{\raggedright\arraybackslash}p{4.6cm} X C{1.9cm}}
\rowcolor{guinda}
\encabezado{Función} & \encabezado{Descripción} & \encabezado{Sección} \\
Gestión de materias & Registro, organización, ocultamiento y eliminación de las materias que imparte, con confirmación en las acciones destructivas. & \ref{subsec:agregar} \\
Carga de presentaciones & Archivos \texttt{.pptx} (máximo 30~MB) agrupados por materia, con miniatura y soporte de varias versiones. & \ref{subsec:cargar} \\
Análisis con IA & Revisión diapositiva por diapositiva con cuatro métricas de calidad, tiempo sugerido de exposición y recomendaciones de escritura. & \ref{subsec:analisis} \\
Versión optimizada & Generación automática de una nueva versión del archivo con las mejoras aplicadas. & \ref{subsec:generar} \\
Cobertura del temario & Comparación de los encabezados de las diapositivas contra los temas oficiales de cada materia. & \ref{subsec:temario} \\
Sesión en vivo & Proyección al grupo con panel de apoyo para el profesor: cronómetros, preguntas y datos complementarios. & \ref{subsec:clase} \\
Informe de rendimiento & Tiempo planeado contra tiempo real por diapositiva y por grupo escolar. & \ref{subsec:rendimiento} \\
Historial de versiones & Registro de cada versión cargada o generada, con su evaluación y recomendaciones. & \ref{subsec:historial} \\
Ocultar y eliminar & Retiro temporal o permanente de presentaciones y materias. & \ref{subsec:ocultar} \\
\end{tabularx}
\end{center}

\subsection{¿Cómo evalúa la calidad de una presentación?}
\label{subsec:metricas}

El análisis revisa cada diapositiva en cuatro aspectos:

\begin{center}
\small
\begin{tabularx}{\linewidth}{>{\raggedright\arraybackslash}p{4.9cm} X >{\raggedright\arraybackslash}p{3.4cm}}
\rowcolor{guinda}
\encabezado{Aspecto} & \encabezado{Qué revisa} & \encabezado{Etiquetas posibles} \\
Complejidad de lectura & Qué tan fácil de leer y entender es el texto de la diapositiva. & Bien / Mejorar \\
Cantidad de palabras & Si el texto es suficiente para explicar el tema sin sobrecargar la diapositiva. & Bien / Mejorar \\
Claridad del encabezado & Si el título refleja realmente el contenido de la diapositiva. & Bien / Mejorar \\
Flujo narrativo & Si existe continuidad entre una diapositiva y la siguiente. & Bien / Revisar \\
\end{tabularx}
\end{center}

El desempeño general de la diapositiva y de la presentación completa se muestra con
etiquetas de calidad: \ui{Excelente}, \ui{Bueno}, \ui{Regular} y \ui{Deficiente}.
Cuando una diapositiva no requiere cambios, se muestra \ui{Óptima}. La aplicación
muestra únicamente estas etiquetas; el detalle del proceso de evaluación no se
muestra al usuario.

\begin{nota}[Tiempo sugerido]
Junto con la calidad, la inteligencia artificial calcula cuánto tiempo conviene
exponer cada diapositiva. Ese tiempo se usa como referencia durante la clase y como
base del informe de rendimiento (sección~\ref{subsec:rendimiento}).
\end{nota}

\subsection{Requisitos del equipo}
\label{subsec:requisitos}

\begin{center}
\small
\begin{tabularx}{\linewidth}{>{\raggedright\arraybackslash}p{5.2cm} X}
\rowcolor{guinda}
\encabezado{Requisito} & \encabezado{Detalle} \\
Equipo Windows & Windows 10 o 11, 64 bits. \\
Conexión a internet & Necesaria durante el análisis; los resultados dependen del servicio de Gemini. \\
Clave de Google Gemini & Se obtiene una vez y se registra desde la propia aplicación (sección~\ref{subsec:configurar}). \\
Proyector o pantalla adicional & Requerido únicamente para la clase en vivo (sección~\ref{subsec:clase}). \\
Presentaciones \texttt{.pptx} & Hasta 30~MB por archivo. \\
\end{tabularx}
\end{center}

\subsection{Convenciones del manual}
\label{subsec:convenciones}

\begin{itemize}
  \item \btn{Botón} es un botón de la aplicación; \ui{texto} es un elemento o mensaje
        que aparece en pantalla; \tecla{tecla} es una tecla del teclado.
  \item Cada paso numerado va seguido de la captura de pantalla correspondiente.
  \item Las referencias como \mbox{la sección~\ref{subsec:cargar}} llevan a la parte
        del manual donde se explica el tema.
\end{itemize}

% ============================================================
%  2. MANUAL DE INSTALACIÓN
% ============================================================
\section{Manual de instalación}
\label{sec:instalacion}

La aplicación se entrega como un archivo ejecutable de Windows. No requiere instalar
programas adicionales ni escribir instrucciones en el equipo: basta con guardar el
archivo y abrirlo. El único requisito externo es la clave de Google Gemini, que se
obtiene desde el navegador y se registra desde la propia aplicación.

\subsection{Requisitos previos}
\label{subsec:previos}

\begin{itemize}
  \item Equipo con Windows 10 o 11, 64 bits.
  \item Conexión a internet.
  \item Clave de Google Gemini (sección~\ref{subsec:clave}).
  \item Proyector o pantalla adicional, solo si va a dar la clase en vivo
        (sección~\ref{subsec:clase}).
\end{itemize}

\subsection{Obtener la aplicación}
\label{subsec:obtener}

\begin{pasos}
  \item Solicite el archivo \texttt{AI Presentation Analyzer.exe} al equipo del
        Trabajo Terminal.
  \item Guárdelo en una carpeta del equipo donde no se mezcle con otros archivos,
        por ejemplo \ui{Documentos\textbackslash AI Presentation Analyzer}.
        La aplicación creará ahí mismo sus carpetas de datos la primera vez que
        registre material.
\end{pasos}

\begin{nota}
No es necesario instalar nada más: el archivo ejecutable contiene todo lo que la
aplicación necesita para funcionar.
\end{nota}

\subsection{Iniciar la aplicación}
\label{subsec:iniciar}

\begin{pasos}
  \item Abra con doble clic el archivo \texttt{AI Presentation Analyzer.exe}.
  \item La primera apertura puede tardar algunos segundos mientras la aplicación
        prepara sus datos. Después se muestra la pantalla principal
        (sección~\ref{subsec:principal}).
  \item Para cerrarla, presione el botón \ui{$\times$} de la ventana.
\end{pasos}

\begin{nota}
Si la aplicación está proyectando una clase, impedirá cerrarse hasta que la
sesión termine (sección~\ref{subsec:clase}).
\end{nota}

\subsection{Obtener la clave de Google Gemini}
\label{subsec:clave}

El análisis de las presentaciones se realiza con la inteligencia artificial de
Google Gemini. Para usarla se necesita una clave personal (\ui{API Key}), que se
obtiene una sola vez desde el navegador:

\begin{pasos}
  \item Entre a \url{https://aistudio.google.com} con su cuenta de Google.
\end{pasos}

\pendiente[4.2cm]{api-01-ai-studio}{Pantalla principal de Google AI Studio}

\begin{pasos}\setcounter{pasosi}{1}
  \item Presione la opción \ui{Get API key} y después \ui{Create API key}.
\end{pasos}

\pendiente[4.2cm]{api-02-create}{Botón Get API key y Create API key}

\begin{pasos}\setcounter{pasosi}{2}
  \item Copie la clave generada con el botón de copia y consérvela; la usará en la
        sección~\ref{subsec:configurar}.
\end{pasos}

\pendiente[4.2cm]{api-03-copiar}{Clave generada y botón de copia}

\begin{importante}
La clave es personal e intransferible. No la comparta ni la publique: quien la
tenga puede consumir el servicio de Gemini a su nombre.
\end{importante}

\subsection{Configurar la clave y el modelo en la aplicación}
\label{subsec:configurar}

La clave se registra \textbf{únicamente desde la aplicación}; no hay archivos que
editar ni configuraciones adicionales.

\begin{pasos}
  \item Inicie la aplicación (sección~\ref{subsec:iniciar}) y presione el botón
        del engranaje \ui{⚙} en la esquina superior derecha.
\imagen[width=0.5\linewidth]{p-engranaje}
  \item Se abre la ventana \ui{Configuración de IA}. Pegue su clave en el campo
        \ui{API Key de Gemini}. Con el botón \ui{Mostrar / ocultar} puede
        comprobar que se copió completo.
\imagen[width=0.62\linewidth]{p-config-ia}
  \item En \ui{Modelo de IA}, elija el modelo que desea usar. La lista se carga
        automáticamente; si no aparece, presione \ui{Actualizar}. El modelo
        recomendado ya viene seleccionado (\ui{gemini-3.5-flash-lite}).
\imagen[width=0.62\linewidth]{p-modelo}
  \item Presione \ui{Guardar}. La aplicación confirma que la configuración quedó
        registrada y queda lista para analizar.
\imagen[width=0.62\linewidth]{p-guardar}
\end{pasos}

\begin{nota}
Si la clave no se registró, el análisis no puede ejecutarse y la aplicación lo
avisa con el mensaje \ui{La API Key no puede estar vacía}.
\end{nota}

% ============================================================
%  3. GUÍA PASO A PASO POR FUNCIÓN
% ============================================================
\section{Guía paso a paso por función}
\label{sec:guia}

\subsection{La pantalla principal}
\label{subsec:principal}

Al abrir la aplicación se muestra la pantalla principal, organizada en cuatro zonas:

\begin{itemize}
  \item \textbf{Panel izquierdo:} lista \ui{Mis Materias} con los botones
        \ui{+ Agregar materia} y \ui{Materias ocultas}.
  \item \textbf{Área central:} las tarjetas de las presentaciones de la materia
        seleccionada y la tarjeta \ui{Subir presentación}.
  \item \textbf{Panel derecho:} \ui{Contenido del curso}, con el temario oficial de
        la materia (sección~\ref{subsec:temario}).
  \item \textbf{Barra superior:} el botón \ui{$\equiv$~Contenido}, que muestra u
        oculta el panel derecho, y el botón del engranaje de configuración
        (sección~\ref{subsec:configurar}).
\end{itemize}

\imagen{p-principal}

\subsection{Configurar la inteligencia artificial}
\label{subsec:configurar-guia}

Antes de analizar una presentación conviene verificar que la clave de Gemini ya
esté registrada. El procedimiento único se describe en la
sección~\ref{subsec:configurar}; desde la guía basta saber que el engranaje
\ui{⚙} de la barra superior abre la ventana \ui{Configuración de IA}.

\subsection{Agregar una materia}
\label{subsec:agregar}

\begin{pasos}
  \item En el panel izquierdo presione \ui{＋ Agregar materia}.
\imagen[width=0.46\linewidth]{p-agregar-btn}
  \item Se abre el catálogo de materias de la ESCOM. Use el buscador, seleccione la
        materia y presione \ui{Agregar materia}.
\imagen[width=0.66\linewidth]{p-agregar-modal}
  \item La materia aparece en la lista \ui{Mis Materias} y queda seleccionada para
        empezar a trabajar con ella.
\imagen[height=8.2cm]{p-mis-materias}
\end{pasos}

\begin{nota}
Puede repetir el procedimiento para registrar todas las materias que imparta. La
aplicación inicia siempre en la primera materia de la lista.
\end{nota}

\subsection{Cargar una presentación}
\label{subsec:cargar}

\begin{pasos}
  \item Seleccione en \ui{Mis Materias} la materia donde va a trabajar.
  \item En el área central presione la tarjeta \ui{Subir presentación}.
\imagen[height=6.2cm]{p-subir}
  \item En el explorador de archivos de Windows elija el archivo \texttt{.pptx} y
        confirme. La aplicación genera la miniatura y la tarjeta de la presentación.
\imagen[height=6.2cm]{p-tarjeta}
\end{pasos}

\begin{importante}[Restricciones al cargar]
\begin{itemize}
  \item Solo se aceptan archivos \texttt{.pptx} de hasta \textbf{30~MB}; un archivo
        mayor muestra el mensaje \ui{El archivo es demasiado pesado}.
  \item No puede cargar dos presentaciones con el mismo nombre dentro de la misma
        materia.
  \item El nombre del archivo no debe ser demasiado largo: Windows limita las
        rutas de carpetas y archivos.
\end{itemize}
\end{importante}

\subsection{Ejecutar el análisis con IA}
\label{subsec:analisis}

\begin{pasos}
  \item Sobre la tarjeta de la presentación presione \btnsuave{Ver opciones}.
        Se abre el menú de la tarjeta.
\imagen[height=6.2cm]{p-menu}
  \item Elija \btn{Analizar presentación}. Aparece un anuncio de progreso mientras
        la aplicación extrae el contenido de cada diapositiva y consulta a Gemini;
        el proceso puede tardar algunos minutos. No cierre la ventana.
  \item Al terminar se abre la \textbf{vista de análisis}: el visor de diapositivas
        a la izquierda y, a la derecha, el panel \ui{Análisis de IA} de la
        diapositiva mostrada, con el tiempo sugerido por la IA, el desempeño
        general y las cuatro métricas con sus etiquetas
        (sección~\ref{subsec:metricas}).
\imagen{p-analisis}
  \item Cuando una diapositiva tiene recomendaciones de escritura, el panel muestra
        el botón \btn{Ver sugerencias de reestructuración}. Al presionarlo se
        despliega el contenido optimizado propuesto; presione \btn{Entendido}
        para cerrarlo.
\imagen[width=0.6\linewidth]{p-sugerencias}
  \item Revise las diapositivas con los botones \tecla{<} y \tecla{>} o con las
        teclas \tecla{$\leftarrow$} y \tecla{$\rightarrow$}. Para aplicar las
        mejoras, presione
        \btn{Generar y aplicar cambios}: la aplicación crea la versión optimizada
        y la abre para que la revise.
\imagen[width=0.55\linewidth]{p-generar}
\end{pasos}

\begin{nota}[¿Dónde queda la versión optimizada?]
La versión optimizada se guarda en la carpeta
\ui{Documentos\textbackslash AI Presentation Analyzer\textbackslash <materia>}
con el nombre \ui{<nombre de la presentación>\_v<versión>.pptx}; por ejemplo, la
versión~2 de una presentación llamada \ui{Unidad1} se guardaría como
\ui{Unidad1\_v2.pptx} dentro de la carpeta de su materia. Si desea buscarla
manualmente, ahí es donde se encuentra. El archivo original que cargó nunca se
modifica.
\end{nota}

\subsection{Consultar la cobertura del temario}
\label{subsec:temario}

\begin{pasos}
  \item Seleccione la materia en \ui{Mis Materias}.
  \item Presione el botón \ui{$\equiv$~Contenido} de la barra superior para
        mostrar el panel derecho.
\imagen[width=0.3\linewidth]{p-contenido-btn}
  \item El panel \ui{Contenido del curso} lista las unidades, temas y subtemas
        oficiales de la materia e indica cuáles ya están abarcados por las
        presentaciones cargadas. Al presionar un tema, la aplicación muestra las
        presentaciones que lo cubren.
\imagen[height=9cm]{p-temario}
\end{pasos}

\subsection{Actualizar una presentación}
\label{subsec:actualizar}

Cuando modifique el archivo \texttt{.pptx} de una presentación ya cargada, no hace
falta cargarlo de nuevo:

\begin{pasos}
  \item Sobre la tarjeta de la presentación presione \btnsuave{Ver opciones}
        (sección~\ref{subsec:analisis}) y elija \ui{Actualizar presentación}.
  \item Seleccione el archivo nuevo en el explorador de Windows.
  \item La aplicación registra el archivo como una versión nueva (v2, v3, \dots)
        y conserva las anteriores. El historial se consulta en la
        sección~\ref{subsec:historial}.
\end{pasos}

\begin{nota}
El archivo original de cada versión nunca se sobrescribe: cada actualización
agrega un archivo distinto en la carpeta de la materia.
\end{nota}

\subsection{Presentar la clase en vivo}
\label{subsec:clase}

\begin{importante}[Restricción del proyector]
Para iniciar una clase se necesita un \textbf{proyector o una segunda pantalla
conectados al equipo}. Sin ese dispositivo la aplicación no proyecta la
presentación al grupo.
\end{importante}

\begin{pasos}
  \item Sobre la tarjeta de la presentación presione \btnsuave{Ver opciones} y
        elija \ui{Presentar clase}.
  \item Se abre la \textbf{vista del docente}: el visor de diapositivas y el panel
        \ui{Apoyo de Clase}, con los cronómetros \ui{Total (Sesión)} y
        \ui{Esta slide} (con el tiempo sugerido por la IA como referencia), además
        de las preguntas de evaluación y los datos curiosos o complementarios de la
        diapositiva mostrada.
\imagen{p-docente}
  \item Presione \btn{Presentar}. Se abre el formulario \ui{Iniciar Clase}:
        \begin{itemize}
          \item \ui{Grupo Escolar}: elija un grupo existente (\ui{Seleccionar}) o
                cree uno nuevo (\ui{Crear Nuevo}); el nombre admite hasta
                5~caracteres alfanuméricos, por ejemplo \ui{3CV2}.
          \item \ui{Tiempo Estimado}: entre 1 y 90 minutos.
          \item \ui{Dispositivo de Salida}: el proyector o la pantalla donde se
                proyectará.
        \end{itemize}
        Después presione \btn{Iniciar Proyección}.
\imagen[height=8cm]{p-iniciar-clase}
  \item La presentación se proyecta en la pantalla elegida para el grupo, mientras
        en la pantalla del docente sigue visible el \ui{Apoyo de Clase} con los
        cronómetros en marcha y el botón \btn{Detener}.
\imagen[width=0.55\linewidth]{p-proyeccion}
\imagen{p-docente-activa}
  \item Controle la clase con el teclado o un presentador remoto:
\end{pasos}

\begin{center}
\small
\begin{tabularx}{0.86\linewidth}{>{\centering\arraybackslash}p{3.4cm} X}
\rowcolor{guinda}
\encabezado{Tecla} & \encabezado{Acción} \\
\tecla{←} & Diapositiva anterior \\
\tecla{→} & Diapositiva siguiente \\
\tecla{Espacio} o \tecla{Enter} & Diapositiva siguiente \\
\tecla{Re Pág} / \tecla{Av Pág} & Anterior / siguiente (botones de un clicker) \\
\end{tabularx}
\end{center}

\begin{pasos}\setcounter{pasosi}{6}
  \item Al terminar la exposición presione \btn{Detener}: la proyección se cierra
        y los tiempos de la sesión quedan registrados para el informe de
        rendimiento (sección~\ref{subsec:rendimiento}).
  \item Para salir de la vista del docente presione
        \ui{$\leftarrow$~Volver a inicio} en la barra superior.
\end{pasos}

\begin{nota}[Reanudar una clase]
Al iniciar una nueva sesión con un grupo que ya tiene sesiones registradas, la
clase retoma en la última diapositiva que el grupo vio, para continuar donde se
quedó.
\end{nota}

\begin{importante}
Si el proyector se desconecta durante la clase, la presentación se detiene
automáticamente por seguridad y la aplicación lo avisa. Reconecte el dispositivo
e inicie de nuevo la clase: la sesión retoma en la última diapositiva vista.
\end{importante}

\subsection{Consultar el informe de rendimiento}
\label{subsec:rendimiento}

\begin{importante}[Requisito]
El informe requiere \textbf{al menos una sesión registrada}. Si la presentación
aún no se ha expuesto, la aplicación lo indica al intentar abrirla.
\end{importante}

\begin{pasos}
  \item Sobre la tarjeta de la presentación presione \btnsuave{Ver opciones} y
        elija \ui{Ver rendimiento}.
  \item El informe muestra los indicadores de la sesión: la calificación del
        análisis, el tiempo recomendado por la IA, el tiempo que planeó el docente,
        la duración real de la clase y la desviación porcentual entre lo planeado y
        lo real. Debajo, la gráfica \ui{Tendencia de tiempo por diapositiva}
        compara el tiempo real (línea punteada) contra el recomendado por la IA
        (línea sólida), diapositiva por diapositiva.
\imagen{p-rendimiento}
  \item Cada grupo escolar tiene su propia pestaña; seleccione la pestaña del grupo
        que desee consultar.
\imagen[width=0.55\linewidth]{p-pestanas}
  \item Para ver el detalle, presione \btn{Desviaciones}: se abre el panel
        \ui{Desviaciones de tiempo por diapositiva}, que separa las diapositivas
        que duraron muy poco de las que duraron demasiado, tomando como referencia
        lo planeado por el docente. Presione \btn{Entendido} para cerrar.
\imagen[width=0.6\linewidth]{p-desviaciones}
\end{pasos}

\begin{nota}
Los tiempos se registran al momento de detener la proyección
(sección~\ref{subsec:clase}); si no detuvo la proyección, la sesión no se
refleja en el informe.
\end{nota}

\subsection{Consultar el historial de versiones}
\label{subsec:historial}

\begin{pasos}
  \item Sobre la tarjeta de la presentación presione \btnsuave{Ver opciones} y
        elija \ui{Ver historial}.
  \item El historial lista todas las versiones de la presentación (cargadas o
        generadas) con su fecha, el número de diapositivas y los botones de cada
        versión: \ui{Ver Evaluación} abre la vista de análisis de esa versión y
        \ui{Sugerencias IA} muestra sus recomendaciones de escritura.
\imagen{p-historial}
  \item En la columna de exportación, cada versión muestra su estado: \ui{Óptima}
        cuando no requiere cambios, o el botón \ui{Generar}, que crea la versión
        optimizada y cambia a \ui{Abrir} cuando el archivo ya existe
        (sección~\ref{subsec:analisis}).
\end{pasos}

\begin{nota}
Una versión recién cargada que aún no se analiza aparece como \ui{Sin análisis};
ejecute el análisis para ver su evaluación.
\end{nota}

\subsection{Ocultar y eliminar elementos}
\label{subsec:ocultar}

\subsubsection{Presentaciones}

\begin{pasos}
  \item Sobre la tarjeta de la presentación presione \btnsuave{Ver opciones} y
        elija \ui{Ocultar} o \ui{Eliminar}.
  \item \ui{Ocultar} retira la presentación de la vista pero conserva todos sus
        datos; podrá restaurarla después.
  \item \ui{Eliminar} borra la presentación con todas sus versiones, análisis y
        sesiones. La aplicación pide confirmación: esta acción no se puede
        deshacer.
  \item Las presentaciones ocultas se consultan con el botón \ui{Ocultas} del
        encabezado del área central, que permite \ui{Restaurar} o
        \ui{Eliminar} cada una.
\imagen[width=0.3\linewidth]{p-ocultas-btn}
\end{pasos}

\subsubsection{Materias}

\begin{pasos}
  \item En el panel izquierdo, presione el ícono de la papelera de la materia.
  \item Se abre la ventana \ui{¿Ocultar o eliminar materia?}:
\imagen[width=0.6\linewidth]{p-materia-dialog}
  \item \ui{Ocultar} retira la materia del panel pero conserva todos sus archivos
        y análisis. \ui{Eliminar todo} borra de forma permanente los archivos y el
        historial. \ui{Cancelar} cierra la ventana sin hacer cambios.
  \item Las materias ocultas se consultan con el botón \ui{Materias ocultas} de la
        parte inferior del panel izquierdo; ahí puede restaurarlas. Las materias
        ocultas conservan todos sus datos.
\imagen[width=0.55\linewidth]{p-materias-ocultas}
\end{pasos}

% ============================================================
%  4. PROBLEMAS PRINCIPALES
% ============================================================
\section{Problemas principales}
\label{sec:problemas}

\subsection{Guía rápida de diagnóstico}
\label{subsec:diagnostico}

\begin{center}
\small
\begin{tabularx}{\linewidth}{>{\raggedright\arraybackslash}p{4.3cm}
                              >{\raggedright\arraybackslash}p{4.6cm}
                              >{\raggedright\arraybackslash}X}
\rowcolor{guinda}
\encabezado{Situación} & \encabezado{Causa probable} & \encabezado{Qué hacer} \\
La aplicación no abre & El archivo se movió de lugar o el antivirus lo bloqueó. & Vuelva a ejecutar el archivo \texttt{.exe}; si persiste, reporte al equipo del Trabajo Terminal. \\
\ui{El archivo es demasiado pesado} & El \texttt{.pptx} supera los 30~MB. & Comprima las imágenes de la presentación o divídala y vuelva a cargarla. \\
Aviso de archivo duplicado & Ya existe una presentación con ese nombre en la materia. & Renombre el archivo o use \ui{Actualizar presentación} (sección~\ref{subsec:actualizar}). \\
Nombre demasiado largo & Windows limita la longitud de rutas de archivos. & Use un nombre más corto y vuelva a cargarlo. \\
\ui{La API Key no puede estar vacía} & La clave de Gemini no se registró. & Siga la sección~\ref{subsec:configurar}. \\
El análisis falla o no responde & Falta conexión a internet o la clave no es válida. & Verifique la conexión y la clave (secciones~\ref{subsec:clave} y~\ref{subsec:configurar}) y reintente. \\
No aparece \ui{Analizar presentación} & La presentación ya fue analizada. & El menú muestra \ui{Ver análisis}; úselo para abrir los resultados. \\
No puedo iniciar la proyección & No hay proyector o pantalla adicional conectados. & Conecte el dispositivo y abra de nuevo el formulario \ui{Iniciar Clase} (sección~\ref{subsec:clase}). \\
La presentación se detiene sola & El proyector se desconectó durante la clase. & Reconecte el dispositivo e inicie de nuevo; la sesión retoma en la última diapositiva vista. \\
\ui{Ver rendimiento} no disponible & La presentación no tiene sesiones registradas. & Dé la clase primero; los tiempos se registran al detener la proyección (sección~\ref{subsec:clase}). \\
La aplicación no me deja cerrar & Hay una proyección activa. & Presione \btn{Detener} y después \ui{$\leftarrow$~Volver a inicio}. \\
\end{tabularx}
\end{center}

\subsection{La aplicación no abre}
\label{subsec:no-abre}

Vuelva a ejecutar el archivo \texttt{AI Presentation Analyzer.exe}. Si el problema
continúa, pruebe en otra carpeta del equipo o reporte el caso al equipo del Trabajo
Terminal: puede tratarse de un bloqueo del antivirus que requiera autorizar el
archivo.

\subsection{Problemas al cargar una presentación}
\label{subsec:problemas-carga}

\begin{itemize}
  \item \ui{El archivo es demasiado pesado}: la presentación supera los 30~MB.
        Comprima las imágenes desde PowerPoint o divida el contenido en dos
        presentaciones.
  \item Aviso de duplicado: dentro de una misma materia los nombres de
        presentación son únicos. Renombre el archivo o use
        \ui{Actualizar presentación} (sección~\ref{subsec:actualizar}).
  \item Nombre demasiado largo: los archivos de Windows tienen un límite de
        longitud en la ruta completa; use nombres breves.
  \item La tarjeta se creó pero la miniatura tarda: espere unos segundos y vuelva
        a seleccionar la materia en \ui{Mis Materias}.
\end{itemize}

\subsection{Problemas con la clave de API y el análisis}
\label{subsec:problemas-ia}

\begin{itemize}
  \item \ui{La API Key no puede estar vacía}: registre la clave como se indica en
        la sección~\ref{subsec:configurar}; es el único lugar donde se configura.
  \item El análisis no responde o falla: verifique la conexión a internet y que la
        clave siga siendo válida (sección~\ref{subsec:clave}). Después reintente
        desde \ui{Analizar presentación}.
  \item El proceso tarda varios minutos: es normal en presentaciones largas; no
        cierre la ventana mientras aparece el anuncio de progreso.
  \item La versión aparece como \ui{Sin análisis} en el historial: ejecute el
        análisis de esa versión (sección~\ref{subsec:analisis}).
\end{itemize}

\subsection{Problemas con la sesión en vivo}
\label{subsec:problemas-clase}

\begin{itemize}
  \item No puedo iniciar la proyección: se necesita un proyector o una segunda
        pantalla conectados (sección~\ref{subsec:clase}). Conéctelos, abra de
        nuevo \ui{Iniciar Clase} y elija el dispositivo en
        \ui{Dispositivo de Salida}.
  \item La presentación se detiene de repente: el proyector se desconectó. La
        aplicación lo avisa y detiene la sesión por seguridad. Reconecte el
        dispositivo e inicie de nuevo la clase con el mismo grupo: retoma en la
        última diapositiva vista.
  \item Los tiempos no aparecen en el informe: los tiempos se registran al presionar
        \btn{Detener}. Si la clase terminó sin detener la proyección, esa sesión no
        se refleja en el informe de rendimiento.
\end{itemize}

\subsection{Problemas con el informe y el historial}
\label{subsec:problemas-informe}

\begin{itemize}
  \item \ui{Ver rendimiento} lo indica como requisito: necesita al menos una
        sesión registrada para esa presentación.
  \item El botón de exportación muestra \ui{Óptima}: la versión ya está en nivel
        óptimo y no requiere reestructuración.
  \item El botón muestra \ui{Abrir} en lugar de \ui{Generar}: la versión optimizada
        ya fue generada antes; la aplicación la abre directamente. El archivo se
        conserva en \ui{Documentos\textbackslash AI Presentation Analyzer}
        (sección~\ref{subsec:analisis}).
\end{itemize}

\subsection{Recomendaciones finales}
\label{subsec:recomendaciones}

\begin{itemize}
  \item Registre la clave de Gemini una sola vez y verifíquela antes de la primera
        clase con análisis.
  \item Detenga siempre la proyección con \btn{Detener} al terminar la clase, para
        que los tiempos queden registrados en el informe.
  \item Use \ui{Actualizar presentación} para nuevas versiones del mismo material
        en lugar de cargar archivos duplicados.
  \item Prefiera \ui{Ocultar} sobre \ui{Eliminar} cuando aún pueda necesitar el
        material: lo oculto se restaura, lo eliminado no.
  \item Conserve el archivo \texttt{.exe} en su carpeta: de ahí la aplicación
        maneja sus datos.
\end{itemize}

\begin{importante}[Soporte]
Si presenta un problema que no se describe en este manual, reporte al equipo del
Trabajo Terminal 2026-B127 indicando lo que estaba haciendo y el mensaje que
apareció en pantalla.
\end{importante}

\end{document}
